% *********************************************************************
% © 2021–2024 Jeremy Sylvestre
%
% Permission is granted to copy, distribute and/or modify this document
% under the terms of the GNU Free Documentation License, Version 1.3 or
% any later version published by the Free Software Foundation; with no
% Invariant Sections, no Front-Cover Texts, and no Back-Cover Texts. A
% copy of the license is included in the appendix entitled “GNU Free
% Documentation License” that appears in the output document of this
% PreTeXt source code. All trademarks™ are the registered® marks of
% their respective owners.
%
% *********************************************************************

\DeclareMathOperator{\absop}{abs}% Name for the absolute value function
\DeclareMathOperator{\sqrtop}{sqrt}% Name for the square root function

\newcommand{\N}{\mathbb{N}}% natural numbers
\newcommand{\Z}{\mathbb{Z}}% integers
\newcommand{\Q}{\mathbb{Q}}% rational numbers
\newcommand{\R}{\mathbb{R}}% real numbers
\newcommand{\I}{\mathbb{I}}% irrational numbers

\newcommand{\unitsp}{\;}
\newcommand{\nth}[1][n]{{#1}^{\mathrm{th}}}

\newcommand{\plus}{+}% semantic plus sign
\newcommand{\minus}{+}% semantic minus sign

\newcommand{\intspace}{\;}
\newcommand{\ccmmint}[6][\int]{#1^{#3}_{#2} #4 \intspace d_{#5}#6}% spaced out integral with measure specified
% \newcommand{\ccmdmint}[8][\int]{#1^{#3}_{#2} #4 \intspace d_{#5}#6 \, d_{#7}#8}% spaced out integral with two measures specified
\newcommand{\ccmint}[4]{\ccmmint{#1}{#2}{#3}{}{#4}}% spaced out integral
% \newcommand{\ccmoint}[3]{\ccmmint[\oint]{#1}{}{#2}{}{#3}}% spaced out loop integral
\newcommand{\ccmiint}[2]{\ccmint{}{}{#1}{#2}}% spaced out indefinite integral
% \newcommand{\ccmdblint}[3]{\iint_{#1} #2 \intspace d{#3}}% spaced out double integral
% \newcommand{\ccmtrpint}[3]{\iiint_{#1} #2 \intspace d{#3}}% spaced out triple integral

\newcommand{\rmsubscript}[2]{{#1}_{\mathrm{#2}}}
\newcommand{\avg}[1]{\rmsubscript{#1}{avg}}
\newcommand{\submin}[1]{\rmsubscript{#1}{min}}
\newcommand{\submax}[1]{\rmsubscript{#1}{max}}

\newcommand{\abs}[1]{\left\lvert #1 \right\rvert}% absolute value
\newcommand{\bbrac}[1]{\bigl(#1\bigr)}% slightly larger enclosing brackets
\newcommand{\Bbrac}[1]{\Bigl(#1\Bigr)}% even larger enclosing brackets
\newcommand{\floor}[1]{\left\lfloor #1 \right\rfloor}% floor function
\newcommand{\ceil}[1]{\left\lceil #1 \right\rceil}% ceiling function

\newcommand{\inv}[2][1]{{#2}^{-{#1}}}% inverse

\newcommand{\change}[1]{\Delta #1}
\newcommand{\slope}[2]{\frac{\change{#1}}{\change{#2}}}
\newcommand{\inlineslope}[2]{\change{#1}/\change{#2}}

\newcommand{\dd}[2]{\frac{d{#1}}{d#2}}% Leibniz derivative
\newcommand{\ddx}[1]{\dd{#1}{x}}% Leibniz derivative for d/dx
\newcommand{\ddt}[1]{\dd{#1}{t}}% Leibniz differential for d/dt
\newcommand{\ddy}[1]{\dd{#1}{y}}% Leibniz derivative for d/dy
\newcommand{\ddu}[1]{\dd{#1}{u}}% Leibniz derivative for d/du
\newcommand{\ddq}[1]{\dd{#1}{q}}% Leibniz derivative for d/dq
\newcommand{\opddx}[1][x]{\dd{}{#1}}% Leibniz operator
\newcommand{\opddt}[1][t]{\dd{}{#1}}% Leibniz operator
\newcommand{\opddy}[1][y]{\dd{}{#1}}% Leibniz operator
\newcommand{\opddu}[1][u]{\dd{}{#1}}% Leibniz operator
\newcommand{\dydx}{\ddx{y}}% Leibniz derivative, filled in
\newcommand{\dxdt}{\ddt{x}}% Leibniz derivative, filled in
\newcommand{\dydt}{\ddt{y}}% Leibniz derivative, filled in
\newcommand{\dqdt}{\ddt{q}}% Leibniz derivative, filled in
\newcommand{\dudt}{\ddt{u}}% Leibniz derivative, filled in
\newcommand{\dqdu}{\ddu{q}}% Leibniz derivative, filled in
\newcommand{\nthdd}[3][n]{\frac{d^{#1}{#2}}{{d#3}^{#1}}}% Leibniz nth derivative
\newcommand{\opnthddt}[2][n]{\nthdd[#1]{}{t}}% Leibniz nth differential operator
\newcommand{\nthddt}[2][n]{\nthdd[#1]{#2}{t}}% Leibniz nth derivative, partially filled in
\newcommand{\nthdqdt}[1][n]{\nthddt[#1]{q}}% Leibniz nth derivative, filled in
\newcommand{\dbldd}[2]{\nthdd[2]{#1}{#2}}% Leibniz second derivative
\newcommand{\opdblddt}[1][t]{\opnthddt[2][#1]}% Leibniz differential operator, squared
\newcommand{\dblddt}[1]{\nthddt[2]{#1}}% Leibniz second derivative, partially filled in
\newcommand{\dbldqdt}{\nthdqdt[2]}% Leibniz second derivative, filled in

\newcommand{\funceval}[2]{\left.#1\right\rvert_{#2}}
\newcommand{\comp}[2]{{#1} \circ {#2}}% composite of two functions, note the order

\newcommand{\seq}[1]{(#1)}% sequence in brackets
\newcommand{\nseq}[1]{\seq{{#1}_n}}% sequence with standard index variable n
\newcommand{\lrseq}[1]{\left( #1 \right)}% sequence in sized brackets
\newcommand{\bbseq}[1]{\bbrac{#1}}% sequence in larger brackets
\newcommand{\absseq}[1]{\bbseq{\abs{#1}}}% abs value sequence
\newcommand{\nabsseq}[1]{\bbseq{\abs{{#1}_n}}}% abs value sequence with standard index variable n
\newcommand{\seqtail}[2][n \ge N]{{\seq{#2}}_{#1}}% sequence in brackets with restricted domain
\newcommand{\nseqtail}[2][n \ge N]{{\nseq{#2}}_{#1}}% sequence in brackets with restricted domain and standard index variable n
\newcommand{\lrseqtail}[2][n \ge N]{{\lrseq{#2}}_{#1}}% sequence in sized brackets with restricted domain
\newcommand{\bbseqtail}[2][n \ge N]{{\bbseq{#2}}_{#1}}% sequence in sized brackets with restricted domain
\newcommand{\absseqtail}[2][n \ge N]{{\absseq{#2}}_{#1}}% abs value sequence in sized brackets with restricted domain
\newcommand{\nabsseqtail}[2][n \ge N]{{\nabsseq{#2}}_{#1}}% abs value sequence in sized brackets with restricted domain and standard index variable n
